% Complete Section 3.3, original PDF 33-37 / printed 23-27.
% block: adjoint-heading page:33 type:section
\subsection{{The adjoint of an unbounded operator}}\label{sec:unbounded-adjoint}

% block: adjoint-functional-intro page:33 type:ordinary-paragraph
{Now suppose that $T$ is any densely defined linear operator from $X$ to $Y$. Then we may define an unbounded linear operator from $Y$ to $X$ called the \emph{adjoint} of $T$ and denoted by $T^*$. To this end we consider, for any $w\in Y$, the linear mapping $D(T)\to\mathbb R$ given by}

% block: adjoint-functional-equation page:33 type:equation
\begin{equation}\label{eq:adjoint-defining-functional}
v\mapsto\langle w,Tv\rangle_Y,\qquad v\in D(T).
\end{equation}

% block: adjoint-domain-definition page:33 type:ordinary-paragraph
{The domain $D(T^*)$ of the adjoint is defined as the set of $w\in Y$ such that the linear functional \eqref{eq:adjoint-defining-functional} is bounded in $X$ norm, i.e., for which there exists a constant $c_w$ such that}
\[
|\langle w,Tv\rangle_Y|\leq c_w\|v\|_X,\qquad v\in D(T).
\]

% block: adjoint-riesz-representation page:33 type:ordinary-paragraph
{For any such $w\in D(T^*)$ there is a unique extension of the map \eqref{eq:adjoint-defining-functional} to a bounded linear functional from all of $X$ to $\mathbb R$. Moreover, by the Riesz representation theorem, there then exists a unique element $x\in X$ such that}
\begin{equation}\label{eq:adjoint-riesz-identity}
\langle w,Tv\rangle_Y=\langle x,v\rangle_X,\qquad v\in D(T),
\end{equation}

% block: adjoint-identity page:33 type:ordinary-paragraph
{and we set $T^*w=x$. In short,}
\[
\langle T^*w,v\rangle_X=\langle w,Tv\rangle_Y,\qquad v\in D(T),
\]
{and $w\in D(T^*)$ exactly when there exists an element of $X$ (necessarily unique) which, when taken for $T^*w$, fulfills this equation.}

% block: rotated-graph-definition page:34 type:ordinary-paragraph
{Many properties of the adjoint follow from a simple relation between the graph of $T$ and the graph of its adjoint $T^*$. Since the latter lives in $Y\times X$ rather than $X\times Y$, we define a variant of the graph of $T^*$, the \emph{rotated graph}, by}
\[
\widetilde\Gamma(T^*)=\{\,(-x,y)\mid(y,x)\in\Gamma(T^*)\,\}=\{\,(-T^*w,w)\mid w\in D(T^*)\,\}\subset X\times Y,
\]
{where the minus sign has been inserted to make the relationship more evident. In fact, the rotated graph of $T^*$ is simply $\Gamma(T)^\perp$, the orthogonal complement in $X\times Y$ of the graph $\Gamma(T)$.}

% block: adjoint-graph-proposition page:34 type:proposition
\Needspace{6\baselineskip}
\begin{Proposition}\label{prop:adjoint-rotated-graph}
{Let $T$ be a densely defined unbounded operator. Then}
\begin{equation}\label{eq:adjoint-rotated-graph}
\Gamma(T)^\perp=\widetilde\Gamma(T^*),
\end{equation}
\begin{equation}\label{eq:graph-closure-rotated-complement}
\overline{\Gamma(T)}=\widetilde\Gamma(T^*)^\perp.
\end{equation}
\end{Proposition}

% block: adjoint-graph-proof page:34 type:proof
\begin{Proof}
{An element $(x,w)$ of $X\times Y$ belongs to $\Gamma(T)^\perp$ if and only if}
\[
\langle x,v\rangle_X+\langle w,Tv\rangle_Y=0,\qquad v\in D(T),
\]
{which is exactly the condition \eqref{eq:adjoint-riesz-identity} for $w$ to belong to $D(T^*)$ with $T^*w=-x$, i.e., for $(x,w)$ to belong to $\widetilde\Gamma(T^*)$. This proves \eqref{eq:adjoint-rotated-graph}. Taking orthogonal complements then gives \eqref{eq:graph-closure-rotated-complement}.}
\end{Proof}

% block: adjoint-closed-consequence page:34 type:ordinary-paragraph
{As an immediate application, we see that the adjoint $T^*$ of a densely defined unbounded operator $T$ is always closed, whether $T$ is closed or not. Moreover, the original operator $T$ is closed if and only if $\Gamma(T)=\widetilde\Gamma(T^*)^\perp$, since $\Gamma(T)=\overline{\Gamma(T)}$ exactly when $T$ is closed.}

% block: adjoint-dense-intro page:34 type:ordinary-paragraph
{As another application of Proposition~\ref{prop:adjoint-rotated-graph}, we show that the adjoint operator of a \emph{closed} densely defined operator is itself densely defined.}

% block: adjoint-densely-defined page:34 type:proposition
\Needspace{6\baselineskip}
\begin{Proposition}\label{prop:closed-adjoint-densely-defined}
{If $T$ is a closed densely defined unbounded operator from $X$ to $Y$, then $T^*$ is a closed densely defined operator from $Y$ to $X$.}
\end{Proposition}

% block: adjoint-dense-proof page:34 type:proof
\begin{Proof}
{To show that $T^*$ is densely defined, we must prove that any $y\in Y$ which is orthogonal to $D(T^*)$ must vanish. Now}
\[
y\perp D(T^*)\iff(0,y)\in\widetilde\Gamma(T^*)^\perp=\Gamma(T),
\]
{with the last equality coming from Proposition~\ref{prop:adjoint-rotated-graph} and the assumption that $T$ is closed. Thus, if $y\perp D(T^*)$, then $y=T0$ and so indeed vanishes.}
\end{Proof}

% block: double-adjoint-intro page:34 type:ordinary-paragraph
{In view of this result, when $T$ is a closed densely defined operator, then so is $T^*$, and we may define $T^{**}$, which is again a closed densely defined operator $X\to Y$. In fact, $T^{**}$ coincides with $T$.}

% block: double-adjoint-proposition page:34 type:proposition
\Needspace{6\baselineskip}
\begin{Proposition}\label{prop:double-adjoint-identity}
{If $T$ is a closed densely defined unbounded operator from $X$ to $Y$, then $T^{**}=T$.}
\end{Proposition}

% block: double-adjoint-proof page:34 type:proof
\begin{Proof}
{Since $T$ is closed, $\Gamma(T)=\widetilde\Gamma(T^*)^\perp$ and, similarly, $\widetilde\Gamma(T^*)=\Gamma(T^{**})^\perp$. Combining these we get $\Gamma(T)=\Gamma(T^{**})^{\perp\perp}$, which equals $\Gamma(T^{**})$, since $T^{**}$ is closed. Since $T$ and $T^{**}$ have the same graphs, they are equal.}
\end{Proof}

% block: null-range-closed-intro page:35 type:ordinary-paragraph
{If $T$ is a closed operator, then its null space $\mathcal N(T)$ is a closed subspace of $X$. This follows directly from the definition of a closed operator. However, the range of a closed linear operator may well not be closed, even in the case of a bounded linear operator (consider, for example, the inclusion of $H^1$ into $L^2$). The following theorem summarizes the relationship between the null spaces and the ranges of $T$ and $T^*$. Its proof is another application of Proposition~\ref{prop:adjoint-rotated-graph}.}

% block: null-range-orthogonality page:35 type:theorem
\Needspace{6\baselineskip}
\begin{Theorem}\label{thm:adjoint-kernel-range-orthogonality}
{Let $T$ be a closed densely defined operator $X$ to $Y$. Then}
\[
\begin{gathered}
\mathcal R(T)^\perp=\mathcal N(T^*),\quad\mathcal N(T)^\perp=\overline{\mathcal R(T^*)},\\
\mathcal R(T^*)^\perp=\mathcal N(T),\quad\mathcal N(T^*)^\perp=\overline{\mathcal R(T)}.
\end{gathered}
\]
\end{Theorem}

% block: null-range-orthogonality-proof page:35 type:proof
\begin{Proof}
{We verify the third equation first:}
\[
\begin{aligned}
z\in\mathcal N(T)&\iff(z,0)\in\Gamma(T)=\widetilde\Gamma(T^*)^\perp\\
&\iff\langle z,T^*w\rangle_X=0\ \forall w\in D(T^*)\iff z\in\mathcal R(T^*)^\perp.
\end{aligned}
\]
{The third equation then implies the first equation by replacing $T$ by $T^*$. Taking orthogonal complements of the first and third equations then gives the fourth and second equations, respectively.}
\end{Proof}

% block: closed-range-theorem-intro page:35 type:ordinary-paragraph
{Although the range of a closed linear operator $T$ may not be closed, if it is, then so also is the range of $T^*$ (and vice versa, since we can apply this result with $T$ replaced by $T^*$). This is a consequence of Banach's closed range theorem. A proof may be found in the more general setting of closed operators on Banach spaces, for example, in [\BookNumber{71}, Chapter VII.5]. For the convenience of the reader we include a proof in the Hilbert space case here. First we state a simple lemma.}

% block: lower-bound-closed-range page:35 type:lemma
\Needspace{6\baselineskip}
\begin{Lemma}\label{lem:operator-lower-bound-closed-range}
{Let $T$ be a closed operator between Hilbert spaces $X$ and $Y$ and suppose that there exists $\gamma>0$ such that}
\[
\|Tv\|_Y\geq\gamma\|v\|_X,\qquad v\in D(T).
\]
{Then $T$ has a closed range.}
\end{Lemma}

% block: lower-bound-closed-range-proof page:35 type:proof
\begin{Proof}
{If $\{v_n\}$ is a sequence in $D(T)$ such that $Tv_n$ converges to some $y$ in $Y$, we must show that $y\in\mathcal R(T)$. The hypothesis implies that the sequence $\{v_n\}$ is Cauchy and so converges to some $x$ in $X$. Since $T$ is closed, $x\in D(T)$ and $y=Tx\in\mathcal R(T)$.}
\end{Proof}

% block: closed-range-theorem page:35 type:theorem
\Needspace{6\baselineskip}
\begin{Theorem}{{closed range theorem}}\label{thm:closed-range}
{If the range of a closed densely defined operator $T$ from $X$ to $Y$ is closed in $Y$, then the range of $T^*$ is closed in $X$.}
\end{Theorem}

% block: closed-range-proof-restriction pages:35,36 type:proof
\begin{Proof}
{Let $W=\mathcal R(T)$, a closed subspace of $Y$, and let $\overline T$ be the range restriction of $T$, i.e., the unbounded operator from $X$ to $W$ with the same domain as $T$ with $\overline Tv=Tv$, $v\in D(T)$. Obviously $\overline T$ is also a closed densely defined operator between Hilbert spaces, and moreover it is a surjection. We shall show that $\mathcal R(\overline T^*)=\mathcal R(T^*)$. This follows from two observations. First, note that $D(\overline T^*)=D(T^*)\cap W$ and $\overline T^*$ is the restriction of $T^*$ to this space. This follows directly from the definition of the adjoint operator and the fact that $\overline Tv=Tv$ for all $v$ in $D(T)=D(\overline T)$. Thus $\mathcal R(\overline T^*)\subset\mathcal R(T^*)$. But $W^\perp=\mathcal N(T^*)$. Therefore, if $y\in D(T^*)$ and $w$ is its orthogonal projection onto $W$, then $w\in D(\overline T^*)$ and $\overline T^*w=T^*y$, showing that $\mathcal R(T^*)\subset\mathcal R(\overline T^*)$.}

% block: closed-range-proof-surjection page:36 type:proof
{In view of this, it suffices to show that if $T$ is a closed densely defined \emph{surjection}, then $\mathcal R(T^*)$ is closed. We know that $D(T)$ is a Hilbert space when endowed with the graph norm. Let $Z$ be the orthogonal complement of $\mathcal N(T)$ in $D(T)$ and let $\widetilde T:Z\to Y$ be the restriction of $T$. Then $\widetilde T$ is a bounded linear operator which is one-to-one and onto. It follows that it is invertible, so there exists a positive constant $c$ such that for any $y\in Y$ there exists a unique $z\in Z$ with $Tz=y$ and $\|z\|_X\leq c\|y\|_Y$. Applying this for $y\in D(T^*)$ we have}
\[
\|y\|_Y^2=\langle y,Tz\rangle_Y=\langle T^*y,z\rangle_X\leq\|T^*y\|_X\|z\|_X\leq c\|y\|_Y\|T^*y\|_X.
\]
{Thus}
\[
\|T^*y\|_X\geq c^{-1}\|y\|_Y,\qquad y\in D(T^*).
\]
{Now we apply the preceding lemma to $T^*$ to conclude that $\mathcal R(T^*)$ is closed.}
\end{Proof}

% block: finite-codimension-intro page:36 type:ordinary-paragraph
{We conclude this section with some useful results for verifying that the range of an operator is closed. First we show that if the range of a closed linear operator is finite codimensional, then it is closed. A proof in the case of bounded operators is given in [\BookNumber{52}, Lemma 19.1.1], and it extends easily to closed operators.}

% block: finite-codimension-closed-range page:36 type:theorem
\Needspace{6\baselineskip}
\begin{Theorem}\label{thm:finite-codimension-closed-range}
{Let $T$ be a closed linear operator between Hilbert spaces $X$ and $Y$. If the quotient space $Y/\mathcal R(T)$ is finite dimensional, then the range of $T$ is closed in $Y$.}
\end{Theorem}

% block: finite-codimension-proof page:36 type:proof
\begin{Proof}
{First consider the case where $T$ is a bounded linear operator from $X$ to $Y$. Restricting $T$ to the orthogonal complement of its null space, we may assume that $T$ is injective. Let $y_1,\ldots,y_n\in Y$ be such that their cosets modulo $\mathcal R(T)$ are a basis for $Y/\mathcal R(T)$. Then every element of $Y$ may be written as $r+c_1y_1+\cdots c_ny_n$ with $r\in\mathcal R(T)$ and $c_i\in\mathbb R$ in a unique way. Thus if we define $\widetilde T:X\times\mathbb R^n\to Y$ by}
\[
\widetilde T(x,c_1,\ldots,c_n)=Tx+c_1y_1+\cdots c_ny_n,
\]
{then $\widetilde T$ is a bounded linear operator which is one-to-one and onto and hence, by Banach's theorem, a homeomorphism. Since $X\times\{0\}$ is closed in $X\times\mathbb R^n$, it follows that $\mathcal R(T)=\widetilde T(X\times\{0\})$ is closed in $Y$.}

% block: finite-codimension-proof-unbounded page:36 type:proof
{In the general case of a closed linear operator $T$ from $X$ to $Y$, we may view $T$ as a bounded linear operator from $D(T)$ to $Y$, and so, from the previous case, if its range is finite codimensional, it is closed.}
\end{Proof}

% block: compactness-criterion-intro page:36 type:ordinary-paragraph
{Our second criterion for closed range is based on compactness.}

% block: compact-domain-closed-range page:36 type:theorem
\Needspace{6\baselineskip}
\begin{Theorem}\label{thm:compact-domain-closed-range}
{Let $T$ be a closed linear operator between Hilbert spaces $X$ and $Y$, and suppose that the inclusion of its domain into $X$ is compact (where the domain is given the graph norm). Then $\mathcal R(T)$ is closed in $Y$.}
\end{Theorem}

% block: compact-domain-proof-boundedness pages:36,37 type:proof
\begin{Proof}
{Given a sequence of elements $x_n\in D(T)$ and $y\in Y$ such that $Tx_n\to y$ in $Y$, we must show that $y\in\mathcal R(T)$. Without loss of generality we may take $x_n\in\mathcal N(T)^\perp$. In this case, the sequence $x_n$ is bounded in $X$, as we now show by contradiction. Indeed, if it were not, we could pass to a subsequence, which we continue to denote $x_n$, with $\|x_n\|_X\to\infty$. Defining $w_n=x_n/\|x_n\|_X\in D(T)\cap\mathcal N(T)^\perp$ we would thus have a sequence with $\|w_n\|_X=1$ and $Tw_n\to0$. Such a sequence is bounded in $D(T)$, and so, by compactness, we could pass to a subsequence which converges in the norm of $X$ to some $w\in X$. Since $T$ is closed, the resulting $w$ would belong to $D(T)$ and satisfy $Tw=0$, i.e., would belong to $\mathcal N(T)$. However, as the limit of the $w_n$, $w$ would also belong to $\mathcal N(T)^\perp$ and be of norm $1$. These conditions are in clear contradiction.}

% block: compact-domain-proof-limit page:37 type:proof
{Thus it suffices to show that for a bounded sequence of elements $x_n\in D(T)$ with $\lim Tx_n=y$ in $Y$, $y$ belongs to $\mathcal R(T)$. We again invoke the compactness hypothesis to replace the sequence with a subsequence for which, in addition, $x_n$ converges to some $x$ in $X$. Since $T$ is closed we have $x\in D(T)$ and $Tx=y$, as desired.}
\end{Proof}
